\chapter{Quasi-Monte Carlo techniques}

A quasi-Monte Carlo\index{quasi-Monte Carlo}
(QMC)\index{QMC|see{quasi-Monte Carlo}} method is an approximation of
the intregal of $f$ over the $d$-dimensional hypercube defined as
\begin{equation}
\int_{[0,1)^d} f(x) dx \approx \frac{1}{n} \sum_{i=1}^{n} f(x_i),
\label{eq:int_moca}
\end{equation}
for carefully chosen points $x_1,\ldots,x_n$ in the unit hypercube
$[0,1)^d$.
In a stochastic framework, we can rewrite \eqref{eq:int_moca} as
\[
E[f(U_1,\ldots,U_d)] \approx \frac{1}{n} \sum_{i=1}^{n} f(x_i).
\]
In contrast with MC method, classical QMC approaches are based on
carefully designed deterministic point sets.
In particular, these point sets can be the first $n$ points of an
infinite sequence, where the points are enumerated in a specific
order.
The construction of the points $x_i$ also depends explicitely on the
problem dimension; in particular, the vectors $x_i$ in $[0,1)^s$
cannot be constructed by taking sets of $d$ consecutive elements from
a scalar sequence.
Moreover, without an upper bound on $s$, QMC methods are inapplicable.

We will enumerate in this chapter the points of a point set $P_n \in
[0,1)^s$ from 0 to $n-1$, for consistency with the usual notation in
the QMC litterature.
Numerous references exist about QMC techniques.
This chapter is partly inspired from Glasserman~\cite{Glas04},
Chapter~8, while we strongly limit the discussion as it would
otherwise leads us outside the scope of this documentation.

\section{Low discrepancy sequences}

As stated before, we want to generate a set of points $x_1$,
$x_2$,\ldots, $x_N$ covering the unit cube $U^s = [0,1)^S$ as
uniformly as possible.
This first requires to formalize the notion of uniformity.
Given a collection of (Lebesgue) measurable subsets of $[0,1)^d$, the
\Index{discrepancy} of the point set $\lbrace x_1, \ldots, x_n \rbrace$
relative to $\mathcal{A}$ is
\[
D(x_1, x_2,\ldots, x_n; \mathcal{A}) = \sup_{A \in \mathcal{A}} \left|
  \frac{\# \lbrace x_i \in A \rbrace}{n} - \mbox{vol}(A) \right|,
\]
where $\# \lbrace x_i \in A \rbrace$ denotes the numbers of $x_i$
contained in $A$ and $\mbox{vol}(A)$ represents the volume (measure)
of $A$.

Taking $\mathcal{A}$ to be the collection of all rectangles in
$[0,1)^d$ of the form
\[
\prod_{j = 1}^d [u_j, v_h),\quad 0 \leq u_j < v_j \leq 1,
\]
we obtain the ordinary (or extreme)discrepancy
\index{discrepancy!ordinary}\index{discrepancy!extreme|see{ordinary discrepancy}} 
$D(x_1,\ldots,x_n)$.
Restricting $\mathcal{A}$ to rectangles of the form
\[
\prod_{j = 1}^d [0,u_j),
\]
defines the star discrepancy\index{discrepancy!star}
$D^*(x_1,\ldots,x_n)$.
From Niederreiter~\cite{Nied92}, Proposition 2.4, we have
\[
D^*(x_1,\ldots,x_n) \leq D(x_1,\ldots,x_n) \leq 2^d D^*(x_1,\ldots,x_n),
\]
so for fixed $d$, the two quantities have the same order of magnitude.

In dimension $d = 1$, Niederreiter also shows that
\[
D^*(x_1,\ldots,x_n) \geq \frac{1}{2n},\quad D(x_1,\ldots,x_n) \geq \frac{1}{n},
\]
and that, in both cases, the minimum is attained by
\[
x_i = \frac{2i-1}{2n},\ i = 1,\ldots,n.
\]
This corresponds to the midpoint rule integration on the unit
interval.

If we now consider a sequence $x_1$, $x_2$,\ldots of points in
$[0,1)$, it can be shown (Niederreiter~\cite{Nied92}, p.~24) that
\[
D(x_1,\ldots,x_n) \geq D^*(x_1,\ldots,x_n) \geq \frac{c \log n}{n},
\]
for infinitely many $n$, with $c$ being a constant.

In dimensions $d \geq 2$, it is widely believed
(Niederreiter~\cite{Nied92}, p.~32) that any point set $x_1$,
$x_2$,\ldots, $x_n$ satisfies
\[
D^*(x_1,\ldots,x_n) \geq \frac{c_d (\log n)^{d-1}}{n},
\]
and the first $n$ elements of any sequence $x_1$, $x_2$,\ldots,
satisfy
\[
D^*(x_1,\ldots,x_n) \geq \frac{c_d' (\log n)^{d}}{n},
\]
for constants $c_d$, $c_d'$ depending only on dimension $d$. It is
therefore customary to use the term ``low-discrepancy'' for methods
that achieve a start discrepancy of $O((\log n)^d/n)$.
The logarithmic term can be absorbed into any power of $n$, allowing
the looser bound $O(1/n^{1-\epsilon})$, for all $\epsilon > 0$.

\subsection{Van der Corput sequences}

While the library does not offer any implementation of Van der Corput
sequences, it is interesting to review them in order to illustrate
some general principles.

Denote a base an integer $b \geq 2$. Every positive integer $k$ has a
unique representation (called its base-$b$ or $b$-ary expandion) as a
linear combination of nonnegative powers of $b$ with coefficients in
$\lbrace 0, 1,\ldots, b-1 \rbrace$.
We can write this as
\[
k = \sum_{j = 0}^\infty a_j(k) b^j,
\]
with all but finitely many of the coefficient $a_j(k)$ equal to zero.
The radical inverse function $\psi_b$ maps each $k$ to a point in
$[0,1)$ by flipping the coefficients of $k$ about the base-$b$
``decimal'' point to get the base-$b$ fraction $.a_0a_1a_2\ldots$.
More precisely,
\[
\psi_b(k) = \sum_{j = 0}^\infty \frac{a_j(k)}{b^{j+1}}.
\]
The base-$b$ Van der Corput sequence is the sequence $0 = \psi_b(0)$,
$\psi_b(1)$, $\psi_b(2)$,\ldots
Theorem~3.6 of Niederreiter~\cite{Nied92} shows that all Van der
Corput sequences are low-discrepancy sequences.

\section{Digital nets}

\begin{align*}
i & = \sum_{\ell = 0}^{k-1} a_{i,\ell} b^{\ell},\\
\begin{pmatrix}
u_{i,j,1} \\ u_{i,j,2} \\ \vdots
\end{pmatrix} & = C_j
\begin{pmatrix}
a_{i,0} \\ a_{i,1} \\ \vdots \\ a_{i,k-1}
\end{pmatrix} \mod b, \\
u_{i,j} & = \sum_{\ell = 1}^{\infty} u_{i,j,\ell}b^{-\ell}, \\
\bu & = (u_{i,1},\ldots,u_{i,t}).
\end{align*}

\begin{verbatim}
void moca_XXX_next(XXX, double *x);
\end{verbatim},
where \verb|XXX| is \verb|faure|, \verb|halton| or \verb|sobol|.

\subsection{\Index{Halton} sequences}

The Halton sequence is a $d$-dimensional generalization of the Van Der
Corput sequence.
Let $(p_1,\ldots,p_d)$ be the $d$ first \index{prime} numbers, then $x_i$ is
defined by:
\[
x_i = (\phi_{p_1}(i),...,\phi_{p_1}(I)),
\]
where $\phi_{p_j}(i)$ is the Van der Corput sequence in base $p_j$.
The Halton sequence satisfies:
\[
D_i^*(x) \leq \frac{1}{i} \prod_{j=1}^d \frac{p_j \log(p_j i)}{\log
  p_j} = O\left( \frac{\log^d(i)}{i} \right),
\]
with a constant $C_d = \prod_{k=1}^{d} \frac{p_k-1}{2 \log p_k}$.
This constant grows to infinity super-exponentially with dimension.

\subsection{Scrambled Halton\index{Halton!scrambled} sequences}

Even though standard Halton sequences perform very well in low
dimensions, correlation problems have been noted between sequences
generated from higher primes. For example if we started with the
primes 17 and 19, the first 17 pairs of points would have perfect
linear correlation. To avoid this, it is common to drop the first 20
entries, or some other predetermined number depending on the primes
chosen. In order to deal with this problem, various other methods have
been proposed; one of the most prominent solutions is the scrambled
Halton sequence, which uses permutations of the coefficients used in
the construction of the standard sequence.

\subsection{\Index{Sobol} sequences}

The Sobol sequence is a $d$-dimensional sequence in base 2 and it is a
$(t,d)$-sequence. It is one of the most used sequences for Quasi-Monte
Carlo simulation. It was first developped by Sobol and it has been
proved to have some additional uniformity property under some
initialization conditions. Its construction is based on
primitive polynomials over $\mathcal{F}_2$ and XOR operations.
 
The generator matrices $C_j$ are upper triangular binary matrices with
1's on the diagonal:
\[
C_j = \begin{pmatrix}
1 & v_{j,1,2} & \ldots & v_{j,1,c} & \ldots \\
0 & 1 & \ldots & v_{j,2,c} & \ldots \\
\vdots & 0 & \ddot & \vdots & \ddots \\
\vdots & \vdots & \ldots & 1 & \ldots
\end{pmatrix}.
\]

To determine the integers $m_{j,c}$, for each $j$, we first select a
primitive polynomial\index{polynomial!primitive} over $\mathcal{F}_2$, say
\[
f_j(z) = z^{d_j} + a_{j,1}z^{d_j-1} + \ldots + a_{j,d_j},
\]
of degree $d_j$, and choose the first $d_j$ integers
$m_{j,1},\ldots,m_{j,d_j}$.

\subsection{Faure sequences}

\section{Other techniques}

\subsection{Modified Latin Hypercube}

\section{Randomized quasi-Monte Carlo}

The construction of a Randomized quasi-Monte Carlo (RQMC)
\index{RQMC|see{quasi-Monte Carlo}}
\index{quasi-Monte Carlo!randomized}
point set starts with a QMC point set $P_n = \lbrace u_0,\ldots, u_{n-1}
\rbrace$, that cover the unit hypercube $[0, 1)^s$ in a very uniform
way, and randomizes $P_n$ so that after the randomization:
\begin{enumerate}[(R.1)]
\item
it retains its high uniformity when taken as a set and
\item
each individual point has the uniform distribution over $[0, 1)^s$.
\end{enumerate}
Let $\tilde{P}_n = \lbrace U_0,\ldots,U_{n-1} \rbrace$ denote the randomized
points.
The estimator of $\mu = E[f(U)]$ based on one copy of the randomization
is
\[
X_{\mbox{rqmc}} = Q_n = \frac{1}{n} \sum_{i = 0}^{n-1} f(U_i).
\]
This randomization is repeated $m$ times, independently, for some positive
integer $m$, with the same $P_n$.
Assume that any points taken from different randomizations are
pairwise independent, in the sense that if $U_i^{(j)}$ is the $i^{\mbox{th}}$
point from the $j^{\mbox{th}}$ randomization, then
\begin{enumerate}[(R.3)]
\item
for every $i_1$, $i_2$ and $j_1 \ne j_2$, $U^{(j_1)}_{i_1}$ and
$U^{(j_2)}_{i_2}$ are independent.
\end{enumerate}

\subsection{Random shifts}

A simple randomization that satisfies Conditions
(R1) to (R3)with a small amount of change in the point set,
for an arbitrary point set Pn, are the random shift modulo 1 and the random
digital shift. They are defined as follows.
For the random shift modulo 1, we simply generate a single point U uniformly
over [0, 1)t and add it to each point of Pn, coordinate-wise, modulo 1. In
the case of a lattice rule, the lattice structure of the points is preserved by the

\section{Transformation to other distributions}

\begin{table}[htbp]
\begin{tabular}{|c|c|c|}
\hline
Description & Symbol & \#parameters \\
\hline
\verb|constant| & \verb|CONSTANT| & 1 \\
\verb|fixed| & \verb|FIXED| & 1 \\
\verb|bernouilli| & \verb|BERNOUILLI| & 1 \\
\verb|uniform| & \verb|UNIFORM| & 2 \\
\verb|normal| & \verb|NORMAL| & 2 \\
\verb|lognormal| & \verb|LOGNORMAL| & 2 \\
\verb|triangular| & \verb|TRIANGULAR| & 2 \\
\verb|truncated_normal| & \verb|TRUNCATED_NORMAL| & 4 \\
\verb|truncated_normal_b| & \verb|TRUNCATED_NORMAL_B| & 2 \\
errc & ERRC \\
\hline
\end{tabular}
\label{tab:moca_distribution}
\caption{Supported distributions}
\end{table}
